% Steering Taxonomy paper -- working draft.
%
% Target venue: ACL/EMNLP Findings (8-page long paper).
% This draft uses the plain `article` class so it compiles anywhere;
% swap to \documentclass[11pt,a4paper]{acl} for the ACL submission.
\documentclass[11pt]{article}
\usepackage[T1]{fontenc}
\usepackage[utf8]{inputenc}
\usepackage[margin=1in]{geometry}
\usepackage{amsmath,amssymb}
\usepackage{booktabs}
\usepackage{graphicx}
\usepackage{hyperref}
\usepackage{url}
\usepackage{xcolor}
\usepackage{microtype}
\usepackage[round]{natbib}

\newcommand{\todo}[1]{\textcolor{red}{\textbf{[TODO: #1]}}}
\newcommand{\result}[1]{\textcolor{blue}{\textbf{[result placeholder: #1]}}}

\title{When Activation Steering Works:\\
       A Geometric Predictor Across 12 Tasks}
\author{Suleman Imdad\\
        Johns Hopkins University\\
        \texttt{simdad1@jh.edu}}
\date{}

\begin{document}
\maketitle

\begin{abstract}
Activation steering has produced striking results on some tasks (refusal,
honesty, sycophancy) and quietly failed on others. We ask: \emph{when}
does it work? We propose a controlled study across 12 tasks spanning
three hypothesized kinds---behavioral, content-specific, and borderline
---and characterize each task's steering direction along two
complementary axes: a \textbf{geometric} axis (split-half cosine
stability and per-pair angular variance of the contrastive direction)
and an \textbf{empirical} axis ($\Delta$ vs baseline and $\Delta$ vs a
random-direction control). Using a single uniform protocol with
identical hyperparameters across tasks, we find that geometric stability
predicts steering success: tasks with a stable, low-variance direction
move under the intervention; tasks with an unstable or high-variance
direction do not, regardless of effect size in the literature. The
prediction holds on the held-out borderline tasks. The taxonomy gives
practitioners a fast, pre-run diagnostic (no held-out generations
needed) for whether a target axis admits steering at all.
\result{rule accuracy} on 12 tasks; \result{borderline accuracy}
on the 4 borderline cases.
\end{abstract}

% ===========================================================================
\section{Introduction}
% ===========================================================================

Activation steering---adding a precomputed direction to a model's
residual stream at inference time---has become a popular tool for
behavioral control of language models~\citep{turner2023activation,
zou2023representation,rimsky2023steering,li2023iti,arditi2024refusal}.
Reported effects span dramatic (40\,pp+ on
context-faithfulness~\citep{yang2026contextfocus} and refusal
suppression~\citep{arditi2024refusal}) to nearly null. The literature
documents the successes; the failures rarely get their own paper. As a
result, we lack a principled answer to a basic question: \emph{given a
candidate behavior, can we predict in advance whether activation
steering will move it?}

We answer ``mostly yes'' with a single diagnostic computed from the
contrastive pair set alone, before any held-out evaluation:
\begin{quote}
\textbf{Geometric stability predicts steering success.} If the per-pair
contrastive directions agree (high split-half cosine, low per-pair
angular variance), the mean direction tends to move the held-out
metric. If they disagree, no single direction reliably moves anything.
\end{quote}

\paragraph{Contributions.}
\begin{itemize}
\item A \textbf{unified protocol} that produces a uniform per-task report
combining \emph{geometric characterization} (split-half cosine,
per-pair variance) and \emph{empirical characterization}
(baseline / steered / random-direction control) of any
contrastive-pair-derived direction.
\item A \textbf{12-task corpus} spanning canonical behavioral targets
(refusal, honesty, sycophancy, sentiment, truthfulness),
content-specific targets (rag-distractor, fact-override,
topic-suppression), and borderline cases (context faithfulness,
persona, politeness, hallucination-grounding).
\item A \textbf{predictive rule} (split-half cosine
$\geq \tau_{\text{stab}}$ AND per-pair variance $\leq \tau_{\text{var}}$
$\Rightarrow$ steering succeeds) tuned on the
behavioral $+$ content-specific subset and validated on the borderline
subset. The rule achieves \result{X / 12} task-level accuracy overall and
\result{Y / 4} on the held-out borderline tasks.
\item Open-source code, task corpus, and per-task reports.
\end{itemize}

% ===========================================================================
\section{Related Work}
% ===========================================================================

\paragraph{Activation steering.} \citet{turner2023activation} introduce
activation addition; \citet{zou2023representation} formalize
representation engineering with framing-pair contrasts; \citet{li2023iti}
introduce Inference-Time Intervention for truthfulness;
\citet{rimsky2023steering} demonstrate Contrastive Activation Addition
across behaviors; \citet{arditi2024refusal} show refusal is mediated by
a single direction; \citet{yang2026contextfocus} push steering to
context-faithfulness on ConFiQA.

\paragraph{When does steering work?} A small but growing set of papers
asks this question piecemeal.
\todo{cite recent surveys + negative-result papers.}
We are not aware of a unified diagnostic combining geometric and
empirical evidence across a broad task corpus.

\paragraph{Geometry of representations.} The split-half cosine
similarity is a standard reliability measure in
cognitive neuroscience and probing literature.
We adapt it to the contrastive-direction setting as a
\emph{pre-evaluation} diagnostic.

% ===========================================================================
\section{Protocol}
% ===========================================================================

A steering task in our corpus exposes:
\begin{enumerate}
\item A set of contrastive pairs $\{(x_i^+, x_i^-)\}_{i=1}^N$ -- two
prompts that differ along the target axis.
\item A set of held-out evaluation examples $\{(p_j, t_j)\}_{j=1}^M$
with a task-specific scoring function
$s(\text{completion}, t_j) \in [0, 1]$.
\end{enumerate}

\paragraph{Direction extraction.}
Run a forward pass on every pair side, take mean-pooled hidden states
at the target layer $\ell$, and form the per-pair direction
$d_i = h(x_i^+; \ell) - h(x_i^-; \ell) \in \mathbb{R}^{d_{\text{model}}}$.
The mean CAA direction is $\bar d = (1/N)\sum_i d_i$, unit-normalized.

\paragraph{Geometric characterization.}
\begin{itemize}
\item \textbf{Split-half cosine} -- random split of the pairs into two
halves, average direction per half, cosine similarity between the two
averages. Repeated 20 times; we report mean $\pm$ std.
High values mean the direction is robust to which half of the pairs you use.
\item \textbf{Per-pair variance} -- $1 - \overline{\cos(d_i, \bar d)}$.
Zero means every per-pair direction is aligned with the mean; one means
they are orthogonal on average. High values mean ``the mean is robust,
but the individual pair directions are scattered around it.''
\end{itemize}
These two metrics are not redundant: a high split-half cosine with high
per-pair variance signals a noisy collection whose mean is stable but
whose individual examples disagree.

\paragraph{Empirical characterization.}
On the held-out evaluation set we measure:
\begin{itemize}
\item \textbf{Baseline} ($b$): no intervention, model's natural output.
\item \textbf{Steered} ($s$): the mean direction is projected \emph{out}
of the residual stream at layer $\ell$ at the final token (scale-free
ablation; this sidesteps brittle additive-scale tuning).
\item \textbf{Random controls} ($r_k, k = 1\ldots K$): same ablation
with a unit-norm random direction. We report mean and std over $K = 5$.
\end{itemize}
We report $\Delta_{\text{base}} = s - b$ and
$\Delta_{\text{rand}} = s - \bar r$.

\paragraph{Predictive rule.}
A task is \emph{predicted to steer} if
\[
\bar{\cos}_{\text{split-half}} \geq \tau_{\text{stab}}
\quad\text{AND}\quad
\sigma_{\text{per-pair}} \leq \tau_{\text{var}}.
\]
A task is \emph{observed to steer} if
$|\Delta_{\text{base}}| \geq \tau_{\text{eff}}$
AND $|\Delta_{\text{rand}}| \geq \tau_{\text{rand}}$.
We tune $\tau$'s on the behavioral $+$ content-specific subset and
report sensitivity analysis.

% ===========================================================================
\section{Task Corpus}
% ===========================================================================

\begin{table}[t]
\centering\footnotesize
\begin{tabular}{l l l}
\toprule
Task & Hypothesized kind & Data source \\
\midrule
refusal              & behavioral       & AdvBench + Alpaca \\
honesty              & behavioral       & TruthfulQA framings \\
sycophancy           & behavioral       & Anthropic MWE \\
sentiment            & behavioral       & IMDB \\
truthfulness         & behavioral       & TruthfulQA \\
\midrule
rag\_distractor       & content-specific & HotpotQA distractor \\
fact\_override        & content-specific & NQ-Swap \\
topic\_suppression    & content-specific & curated \\
\midrule
context\_faithfulness & borderline       & NQ-Swap \\
persona              & borderline       & Alpaca + personas \\
politeness           & borderline       & Alpaca + register \\
hallucination\_grounding & borderline    & TruthfulQA \\
\bottomrule
\end{tabular}
\caption{The 12-task corpus. Behavioral tasks have strong literature priors
for stable single-direction steering; content-specific tasks have a
per-example target that varies by content; borderline tasks plausibly
share a global disposition but vary in surface content.}
\label{tab:corpus}
\end{table}

Table~\ref{tab:corpus} lists the tasks. Each contrastive pair set is
constructed from the data source so that the \emph{difference} between
positive and negative isolates the target axis (e.g.\ for refusal,
positive = harmful prompt, negative = harmless instruction; for
hallucination-grounding, positive = ``I don't know'' continuation,
negative = a confident-but-wrong completion).

% ===========================================================================
\section{Experimental Setup}
% ===========================================================================

\paragraph{Model.}
All experiments use \texttt{meta-llama/Llama-3.2-3B-Instruct}, target
layer $\ell = 7$ of 28 (early-middle), $N = 200$ contrastive pairs,
$M = 100$ held-out evaluation examples, $K = 5$ random controls.
Single fixed seed.

\paragraph{Decoding.}
Greedy, $\leq 64$ new tokens. Steering applies the
ablation at the last input-token position on layer $\ell$'s
residual stream output.

\paragraph{Implementation.}
A \texttt{steering\_taxonomy} Python package; one module per task; one
uniform CLI runs the full corpus.

% ===========================================================================
\section{Results}
% ===========================================================================

\result{Cross-task table (12 rows from analyze\_taxonomy.py).}

\result{Main figure: scatter of per-pair-variance vs split-half-cosine
with effect coloring (Figure A) + bar chart of $|\Delta_{\text{rand}}|$
per task (Figure B). See plot\_taxonomy.py output.}

\result{Predictive-rule accuracy table -- overall, behavioral subset,
content-specific subset, borderline subset.}

\paragraph{Behavioral tasks.} \result{summary of effect sizes on the 5
behavioral tasks.}

\paragraph{Content-specific tasks.} \result{summary -- expected to fail
on the steering effect despite (potentially) high split-half cosine.}

\paragraph{Borderline tasks.} \result{the 4 borderline tasks are the
held-out validation set for the predictive rule.}

% ===========================================================================
\section{Discussion}
% ===========================================================================

\result{Discussion will be written once results land. Key threads to
address:}
\begin{itemize}
\item \textbf{When the rule fails.} Tasks where the rule mispredicts
are the most diagnostic case. Are they ones where the contrastive
construction has a confound?
\item \textbf{Why per-pair variance matters.} A high split-half cosine
with a high per-pair variance is the most interesting geometric case:
the average is stable, but individual examples disagree. Our results
will show whether this signature predicts no-effect or weak effect.
\item \textbf{Practical implications.} The diagnostic only requires
forward passes on the pair set, no held-out generations. A practitioner
can decide whether to invest in evaluation before committing GPU time.
\end{itemize}

% ===========================================================================
\section{Limitations}
% ===========================================================================

Single model (Llama-3.2-3B-Instruct), single layer (7), single seed.
We do not vary the ablation strength because the ablation is scale-free
by construction. The lexicon-based scorers (sentiment, politeness,
hallucination-grounding) underestimate true effect sizes relative to
LLM-as-judge; this hurts the empirical axis but does not bias the
geometric axis, which is computed without generation. Anthropic
model-written-evals (sycophancy) and NQ-Swap (fact-override,
context-faithfulness) inherit their corpora's known biases.

% ===========================================================================
\section{Conclusion}
% ===========================================================================

\result{Conclusion paragraph after results.}

% ===========================================================================
\bibliographystyle{plainnat}
\bibliography{refs}

\end{document}
